\section{Discussion}
\label{sec:discussion}

Consider what it takes to pull a spokesperson off their talking points, and compare it
with what it takes to talk them out of a conviction. Talking points are written by
someone else and carried into the room. They can be quoted back at the speaker or
contradicted by a sharper interlocutor, and when the speaker abandons them, everyone
present can tell the moment it happens. A conviction is nowhere in particular. It shows
up as a tilt in what its owner notices, which objections strike them as serious, and
where they stop conceding. A skilled opponent can strip a speaker of their talking
points in minutes. A conviction yields only to evidence.

Every deployed language model works from talking points. The system prompt is a set of
them, written by the deploying party and carried into a conversation the interlocutor
helps write. Whatever a model is asked to be, its instructions occupy the same context
window its opponent argues in, subject to the same attention, weighed by the same chain
of thought that weighs the opponent's rebuttals. We measured what this arrangement
costs and found the two properties everyone bundles together sitting on opposite sides
of a channel boundary: the talking points are adopted instantly and abandoned half the
time; the compiled tilt is adopted reluctantly and holds against nearly everything
except a good reason.

The result is a strange inversion of the obvious procedure. The obvious way to make a
model believe something is to tell it, and telling works, at first: 98\% adoption. What
telling cannot do is make the belief the model's own, because a told belief remains a
quotable object in the conversation, and anything quotable is editable by whoever
speaks next. The compiled belief is never quoted. It conditions the forward pass
without appearing in it, which means the one place the opponent can reach contains
nothing of it to attack. The model defends the position without reciting the
instruction to defend it.

This is why the finding matters beyond its numbers. Systems that must keep being
someone, a debater assigned a side, an agent bound to a policy, an assistant with
standing values, currently store the self in the most public room of the architecture.
The failures are already in the case law. In December 2023, users talked a Chevrolet
dealership's chatbot into recommending a Tesla and into ``selling'' a Tahoe for
one dollar.\footnote{Chevrolet of Watsonville, December 2023; widely reported,
e.g.\ \emph{Business Insider}, ``A car dealership added an AI chatbot to its site. Then,
all hell broke loose,'' Dec.\ 19, 2023.} In February 2024, a tribunal held Air Canada
liable for a refund policy its chatbot invented under a grieving customer's questioning,
awarding him \$812.02 in damages.\footnote{\emph{Moffatt v.\ Air Canada}, 2024 BCCRT
149.} The remedy on offer everywhere is sharper talking points. Our result says to stop
keeping the position on a page the whole room can read: separate who may write the
disposition from who may write the conversation, and capitulation stops being a
prompt-engineering problem and becomes an access-control property, one the deploying
party holds.

An access-control claim is only as good as its resistance to someone trying to defeat
it, so we tested the compiled channel against a real adversary rather than a single
refusable override. Three multi-turn attack strategies each target a different
channel's easiest point of attack, and \condF{} has the lowest flip rate of any
condition under every one of them.\expref{app:exp-inject} The context channel's apparent robustness
against the earlier, single-shot version of this attack turns out to have been the
belief text sitting unevicted in the model's own context, re-read at the final probe
rather than defended; once the attack denies it that re-reading, its flip rate rises
by 21--46 percentage points. One strategy also reaches \condF{} (29.2\%), but part of
its effect is a generic authority-framing pressure that moves the bare condition too,
rather than a handle found on the compiled channel specifically, a caveat the headline
number should carry along with it.

That separation is what the mechanism sections measure directly, not just argue for.
The two attention channels the write function touches split the labor cleanly:
query modulation carries a content-blind capacity to hold, value modulation carries the
content held (Section~\ref{sec:why-channels-differ}). What carries the effect is which
direction gets written, not how much machinery reads it back out at inference time:
persistence survives truncation to a single direction, and a frozen offset holds as
well as one that responds to the live input\iflongversion{} (Section~\ref{sec:isolation-rank})\fi.
The write has to out-muscle the model's own prior; the hold, once written, competes
with nothing (Section~\ref{sec:dissociation-prior}).

The separation would be worth little if it produced a zealot. It does not. The compiled
model concedes to well-warranted counterevidence at \(2.4\times\) its concession rate to
hollow challenges, and never yields to pressure that offers no evidence at all. That
combination, update for reasons, hold against everything else, is the disposition that
debate-based oversight has assumed its debaters possess and has had no way to install.
It is also the disposition we ask of a good juror or a good scientist.
